\chapter{Diseño e implementación} % Main chapter title

\label{Chapter3} % Change X to a consecutive number; for referencing this chapter elsewhere, use \ref{ChapterX}

% Todos los capítulos deben comenzar con un breve párrafo introductorio que indique cuál es el contenido que se encontrará al leerlo.  La redacción sobre el contenido de la memoria debe hacerse en presente y todo lo referido al proyecto en pasado, siempre de modo impersonal.

\definecolor{mygreen}{rgb}{0,0.6,0}
\definecolor{mygray}{rgb}{0.5,0.5,0.5}
\definecolor{mymauve}{rgb}{0.58,0,0.82}

%%%%%%%%%%%%%%%%%%%%%%%%%%%%%%%%%%%%%%%%%%%%%%%%%%%%%%%%%%%%%%%%%%%%%%%%%%%%%
% parámetros para configurar el formato del código en los entornos lstlisting
%%%%%%%%%%%%%%%%%%%%%%%%%%%%%%%%%%%%%%%%%%%%%%%%%%%%%%%%%%%%%%%%%%%%%%%%%%%%%
\lstset{ %
  backgroundcolor=\color{white},   % choose the background color; you must add \usepackage{color} or \usepackage{xcolor}
  basicstyle=\footnotesize,        % the size of the fonts that are used for the code
  breakatwhitespace=false,         % sets if automatic breaks should only happen at whitespace
  breaklines=true,                 % sets automatic line breaking
  captionpos=b,                    % sets the caption-position to bottom
  commentstyle=\color{mygreen},    % comment style
  deletekeywords={...},            % if you want to delete keywords from the given language
  %escapeinside={\%*}{*)},          % if you want to add LaTeX within your code
  %extendedchars=true,              % lets you use non-ASCII characters; for 8-bits encodings only, does not work with UTF-8
  %frame=single,	                % adds a frame around the code
  keepspaces=true,                 % keeps spaces in text, useful for keeping indentation of code (possibly needs columns=flexible)
  keywordstyle=\color{blue},       % keyword style
  language=Python,                % the language of the code
  %otherkeywords={*,...},           % if you want to add more keywords to the set
  numbers=left,                    % where to put the line-numbers; possible values are (none, left, right)
  numbersep=5pt,                   % how far the line-numbers are from the code
  numberstyle=\tiny\color{mygray}, % the style that is used for the line-numbers
  rulecolor=\color{black},         % if not set, the frame-color may be changed on line-breaks within not-black text (e.g. comments (green here))
  showspaces=false,                % show spaces everywhere adding particular underscores; it overrides 'showstringspaces'
  showstringspaces=false,          % underline spaces within strings only
  showtabs=false,                  % show tabs within strings adding particular underscores
  stepnumber=1,                    % the step between two line-numbers. If it's 1, each line will be numbered
  stringstyle=\color{mymauve},     % string literal style
  tabsize=2,	                   % sets default tabsize to 2 spaces
  title=\lstname,                  % show the filename of files included with \lstinputlisting; also try caption instead of title
  morecomment=[s]{/*}{*/}
}


Este capítulo detalla el diseño y la implementación del sistema predictivo. Se analizan los conjuntos de datos utilizados, la extracción de atributos del audio y el tratamiento de datos según las distribuciones de las variables objetivo. Asimismo, se describen las técnicas de reducción y selección de variables, el proceso de optimización de los modelos y los resultados de la iteración productiva. Finalmente, se presenta la arquitectura desarrollada para el despliegue y la integración del sistema en un entorno operativo.

%----------------------------------------------------------------------------------------
%	SECTION 1
%----------------------------------------------------------------------------------------
\section{Conjuntos de datos}
\label{sec:conjuntos_datos}
% Aquí se describirá la extracción de datos del audio.

\subsection{Análisis de conjuntos de datos públicos}

La definición de la estrategia para la extracción de atributos se fundamenta en una revisión de conjuntos de datos públicos estandarizados en el ámbito MIR. 
El objetivo fue entender cómo repositorios de referencia, como FMA y Million Song Dataset, representan y estructuran los datos musicales, para decidir su viabilidad en el contexto de aplicación de este trabajo. 

\subsubsection{GTZAN}
Este conjunto de datos se considera el equivalente al MNIST \cite{lecun1998gradient} para el análisis de audio \cite{sturm2014state}, dada su adopción histórica como estándar de evaluación en tareas de clasificación de géneros musicales.
\begin{itemize}
    \item Variables predictivas principales: tempo, representaciones chroma, amplitud media cuadrática (RMS), descriptores espectrales y coeficientes MFCC.
    \item Estructura: 1~000 archivos de audio en crudo limitados a 30 segundos, acompañados de metadatos tabulares e imágenes de espectrogramas de Mel en su versión distribuida mediante Kaggle \citep{gtzan_kaggle}. 
    Todas estas representaciones fueron extraídas utilizando la biblioteca Librosa. Las tablas contienen los atributos calculados sobre la pista completa y también sobre divisiones de 3 segundos para incrementar el volumen de datos de entrenamiento. En ambos escenarios se aplica una reducción de dimensionalidad temporal: las variables se calculan inicialmente sobre micro-ventanas de milisegundos y luego se colapsan a valores estáticos mediante estadísticos de agregación, como la media y la varianza.
    \item Variables objetivo principales en Plusyc: se podría relacionar con \textit{suggested genres}. 
    \item Desventajas: provee segmentos cortos de audio, lo que limita la representatividad de la información musical a nivel de canción; y solo abarca 10 géneros musicales, mientras que Plusyc trabaja con más de 500.
\end{itemize}

\subsubsection{FMA estándar}
Se analiza la versión estándar de FMA, diseñada para superar las limitaciones de tamaño y acceso a las colecciones completa y mediana.
\begin{itemize}
    \item Variables predictivas principales: representaciones chroma, Transformada de Q Constante (CQT), tonnetz, MFCC, RMS, descriptores espectrales y tasa de cruces por cero (ZCR).
    \item Estructura: ofrece recortes de audio de 30 segundos para sus subconjuntos estándar, acompañados de matrices de atributos precomputados distribuidos en formato tabular (CSV). Los atributos fueron extraídos utilizando Librosa y posteriormente colapsados mediante estadísticos de agregación, como la media, desviación estándar, asimetría y curtosis, para obtener valores estáticos por pista.
    \item Metadatos contextuales: información tabular adicional que incluye álbum, artista, popularidad, fechas de lanzamiento y ubicación geográfica.
    \item Variables objetivo principales: géneros musicales estructurados en una taxonomía jerárquica de clases y para un subconjunto de aproximadamente 13~000 pistas, proporciona atributos como \textit{danceability, energy, valence, speechiness} y \textit{liveness}.
    \item Desventajas: el análisis está restringido a segmentos de 30 segundos en la versión estándar. Las versiones media y completa exceden las capacidades de procesamiento de la infraestructura actual de la empresa. Además, la incertidumbre sobre las versiones exactas de Librosa utilizadas originalmente dificulta la replicabilidad técnica sobre canciones nuevas. Finalmente, la información etiquetada de 13~000 pistas es mucho menor que la de más de 60~000 de Plusyc.
\end{itemize}

\subsubsection{Million Song Dataset (MSD)}
Desarrollado por la Universidad de Columbia y The Echo Nest \cite{bertin2011million}, este conjunto escaló la investigación de información musical a nivel industrial.
\begin{itemize}
    \item Variables predictivas principales: representaciones de tono (chroma y pitch) y descriptores de timbre, ambos calculados a nivel de segmento.
    \item Estructura: colección de metadatos y descriptores precomputados empaquetados en archivos HDF5.
    \item Método de extracción: a diferencia de la segmentación por ventanas fijas, el MSD utiliza una división adaptativa basada en eventos. Esto fragmenta la canción en segmentos de duración variable. Para cada uno de estos segmentos, se calculan coeficientes de tono y timbre, y se genera una matriz de atributos alineada con la estructura temporal de la pista.
    \item Variables objetivo principales: proporciona atributos musicales fundamentales como \textit{tempo, loudness, key} y \textit{mode}.
    \item Desventajas: los descriptores se generaron con algoritmos propietarios de The Echo Nest, lo que complica procesar la segmentación por eventos en canciones nuevas de forma idéntica para realizar el proceso de inferencia.
\end{itemize}

\subsubsection{AudioSet}
Este conjunto masivo establece un estándar para la clasificación de eventos sonoros generales
\begin{itemize}
    \item Estructura: conjunto de identificadores de videos de YouTube y marcas de tiempo que referencian segmentos de 10 segundos, anotados bajo una ontología de 527 clases.
    \item Variables predictivas principales: el corpus original no extrae descriptores tradicionales. Su distribución oficial provee embeddings precomputados mediante el modelo VGGish \cite{hershey2017cnn}, que sirven como entrada rápida para entrenar clasificadores de terceros sin necesidad de descargar el audio original.
    \item Variables objetivo principales: etiquetas multiclase que identifican la presencia de instrumentos musicales, como tipos de voces (habla, canto) y ambientes acústicos.
    \item Desventajas: el enfoque en sonidos generales limita su precisión para los atributos musicales.
\end{itemize}

% En la tabla \ref{tab:variables_datasets} se presentan algunas de las variables predictivas más relevantes de cada conjunto de datos.

Todos estos conjuntos de datos presentan varias limitaciones frente a la solución propuesta: algunos restringen el análisis a segmentos cortos de audio, mientras que otros requieren tamaños de almacenamiento excesivamente pesados o estructuras de datos difíciles de replicar durante el proceso de inferencia y dependientes de las versiones de las bibliotecas de extracción.

Por otro lado, estas restricciones motivaron que el proceso de extracción se ejecute directamente sobre el catálogo privado para eliminar las limitaciones, y su revisión permitió sentar las bases para determinar cuáles son las variables predictivas en la música, cómo se segmentan las canciones, cómo se reduce la dimensionalidad y para hacer los primeros experimentos predictivos de las variables objetivo requeridas.

Cabe mencionar que, a pesar de que en algunas de las fuentes de datos públicas se incluye información del contexto como la región del artista, el año de lanzamiento, el álbum o las letras, en esta solución se trabaja exclusivamente con información musical que se extrae del audio. 

\subsection{Reducción de dimensionalidad mediante estadísticas}

En el ámbito del procesamiento de señales, la información de un archivo de audio se representa como una serie de tiempo multivariada: una secuencia de puntos de datos \cite{muller2015fundamentals}, en este caso, descriptores numéricos y semánticos, indexados y ordenados cronológicamente.

La estrategia adoptada se inspiró en los conjuntos de datos de referencia, donde la señal se divide en ventanas temporales y la información de cada ventana se colapsa en valores estáticos. Esto permite reducir estas secuencias a un vector por cada canción, buscando optimizar el uso de la memoria y los tiempos de procesamiento.
Esta transformación de la señal de secuencias de números a un vector de estadísticos ocurre de la siguiente manera:

\begin{itemize}
    \item Segmentación temporal: se estableció una ventana de análisis de 10 segundos continuos. Esta decisión se fundamenta en la arquitectura de la red neuronal CNN14, entrenada sobre el conjunto de datos AudioSet, cuyos segmentos originales son de 10 segundos.
    \item Extracción de variables predictivas: para cada segmento de 10 segundos de la canción, el algoritmo usa Librosa y Essentia para extraer el conjunto de variables predictivas principales que proveen las fuentes de datos públicos. También se extrae información global de la canción y algunas de las variables predictivas escalares se calculan de forma global y también de forma segmentada y agregada.
    \item Extracción de embeddings y probabilidades: para esos segmentos de 10 segundos, el algoritmo también extrae un vector de embeddings y las probabilidades de activación para 173 clases de eventos musicales provistas por PANNs inference.
    \item Agrupación estadística: dado que una canción estándar produce una cantidad variable de segmentos de 10 segundos, se aplicó una técnica de colapso estadístico a lo largo del eje temporal, diferenciando el tratamiento según el tipo de variable:
    \begin{itemize}
        \item Para las variables predictivas (DSP), se calcularon siete estadísticos por pista: media, desviación estándar, asimetría, curtosis, y los percentiles 10 (cerca del mínimo), 50 (mediana) y 90 (cerca del máximo para evitar ruido).
        \item Para los embeddings y las probabilidades semánticas de las 173 clases seleccionadas de PANNs, se calcularon cuatro medidas: media, valor máximo, desviación estándar y el percentil 90.
    \end{itemize}
\end{itemize}

Este enfoque matemático permite comprimir el volumen masivo de información de una serie de tiempo en un único vector unidimensional por canción.

Algunos experimentos preliminares sirvieron para validar que esta representación permitía capturar la información esencial de la canción de manera eficiente y compacta.
El algoritmo de extracción, es el resultado de múltiples pruebas con distintas configuraciones de variables, estadísticos y métodos de segmentación. Algunos de los experimentos que no resultaron exitosos comprenden:
pruebas con solo variables predictivas sin embeddings, demuestran que no es posible capturar eventos como el público en vivo o el habla en los primeros intentos de inferencia de \textit{liveness} y \textit{speechiness}. Las pruebas con embeddings pero sin variables predictivas en general no proporcionan la información de \textit{key}, \textit{tempo}, \textit{loudness}, \textit{mood} o \textit{genres}.
Pruebas con distintas ventanas resultaron ser ineficientes en el proceso de extracción de más de 60~000 canciones. Pruebas aplicando todos los estadísticos resultaron en un espacio de variables que consume mucho espacio de almacenamiento y con alta correlación.

\subsection{Conformación del espacio de variables masivo}
% Tras validar la superioridad de la arquitectura mixta, se ejecutó el proceso de extracción sobre el catálogo completo. El procesamiento computacional requirió más de una semana de ejecución ininterrumpida para procesar las canciones completas.

% El algoritmo de extracción definitivo combinó las funciones de Librosa y Essentia con la inferencia de la red neuronal CNN14. El resultado fue un conjunto de datos masivo compuesto por más de 9000 variables predictivas independientes por cada registro de audio. Este extenso espacio de características garantiza la disponibilidad de toda la información acústica y semántica posible para la etapa de experimentación. El alto costo computacional de la extracción justifica esta estrategia de recopilación exhaustiva, eliminando la necesidad de reprocesar el audio y trasladando el desafío técnico hacia la selección rigurosa de las variables más relevantes durante la etapa de modelado, proceso que se detalla en la sección \ref{sec:seleccion_variables} %.


\section{Ejecución de tareas sobre el conjunto de datos}
\label{sec:ejecucion_tareas}
% Aquí se detallarán las distribuciones de las variables objetivo, dificultades encontradas y estrategias de muestreo.


\section{Selección de variables de entrada}
\label{sec:seleccion_variables}
% Aquí se describirán las técnicas de selección y reducción para una arquitectura mixta.


\section{Modelado y optimización}
\label{sec:modelado_optimizacion}
% Aquí se abordará el proceso de modelado y la búsqueda de hiperparámetros.


\section{Iteración productiva}
\label{sec:iteracion_productiva}
% Aquí se presentarán los resultados generales de la iteración.


\section{Arquitectura de despliegue e integración}
\label{sec:arquitectura_despliegue}
% Aquí se documentará la arquitectura para el despliegue e integración en un ambiente productivo.

% La idea de esta sección es resaltar los problemas encontrados, los criterios utilizados y la justificación de las decisiones que se hayan tomado.
% Se puede agregar código o pseudocódigo dentro de un entorno lstlisting con el siguiente código:




\begin{verbatim}
\begin{lstlisting}[caption= "un epígrafe descriptivo"]
	las líneas de código irían aquí...
\end{lstlisting}
\end{verbatim}

A modo de ejemplo, se muestra el fragmento de código \ref{cod:vControl}:

\begin{lstlisting}[label=cod:vControl,caption=Pseudocódigo del lazo principal de control.]  % Start your code-block

#define MAX_SENSOR_NUMBER 3
#define MAX_ALARM_NUMBER  6
#define MAX_ACTUATOR_NUMBER 6

uint32_t sensorValue[MAX_SENSOR_NUMBER];		
FunctionalState alarmControl[MAX_ALARM_NUMBER];	//ENABLE or DISABLE
state_t alarmState[MAX_ALARM_NUMBER];						//ON or OFF
state_t actuatorState[MAX_ACTUATOR_NUMBER];			//ON or OFF

void vControl() {

	initGlobalVariables();
	
	period = 500 ms;
		
	while(1) {

		ticks = xTaskGetTickCount();
		
		updateSensors();
		
		updateAlarms();
		
		controlActuators();
		
		vTaskDelayUntil(&ticks, period);
	}
}
\end{lstlisting}



